% Chapter 0: Introduction and scope
% Written in the voice of the kernel README, oriented toward the blueprint
% audience (learning theorists who do not read Lean 4).

\chapter{Introduction}\label{chap:intro}

This is the mathematical blueprint for a Lean 4 formalization of the
fundamental theorem of statistical learning. It sits between two other
artifacts: the kernel source code at
\url{https://github.com/Zetetic-Dhruv/formal-learning-theory-kernel},
which contains 354 machine-checked theorems in 21{,}728 lines with zero
\texttt{sorry}, and a companion textbook on formal learning theory from
which most of the exposition in Chapters~\ref{ch:pac}
through~\ref{ch:separations} is imported directly. The blueprint's
role is to annotate informal mathematical statements with hyperlinks
into the formal declarations that prove them, so that a reader can
click from a theorem in prose to the corresponding
\texttt{theorem} or \texttt{def} in the Lean kernel and back.

Every statement carrying a \texttt{\textbackslash lean\{\}} annotation
in this document resolves to a proved declaration in the kernel.
The formalization targets Lean 4 \texttt{v4.29.0-rc6} against mathlib4
pinned at commit \texttt{fde0cc5}. The per-module API reference is at
\url{https://zetetic-dhruv.github.io/formal-learning-theory-kernel/}.

\section*{Scope}

Blueprint v1 covers the core story of the kernel:

\begin{itemize}[leftmargin=*,itemsep=3pt]
  \item \textbf{PAC characterization} (Chapter~\ref{ch:pac}): the
    five-way equivalence for PAC learning, with the
    \texttt{NullMeasurableSet} refinement of the standard Borel
    hypothesis used in the literature.
  \item \textbf{Online characterization} (Chapter~\ref{ch:online}):
    the Littlestone characterization theorem with the Standard Optimal
    Algorithm as the constructive witness.
  \item \textbf{Gold-style learning} (Chapter~\ref{ch:gold}): Gold's
    theorem with the diagonalization proof and the mind-change
    characterization with ordinal bounds.
  \item \textbf{Three-paradigm separation}
    (Chapter~\ref{ch:separations}): the 13-edge separation lattice
    between PAC, online, and Gold-style learning.
  \item \textbf{Measurability layer} (Chapter~\ref{chap:measurability}):
    the Borel-parameterized setting, the analytic measurability bridge
    via Choquet capacitability, and the Borel-analytic separation
    theorem. This chapter is net-new: it does not exist in the
    companion textbook because the Borel-analytic separation was
    discovered during the formalization effort.
\end{itemize}

Deferred to v2: compression via approximate minimax, the measurable
batch learner monad, PAC-Bayes bounds, extended criteria (robust PAC,
RKHS), and the Baxter multi-task base case. Compression and the MBL
monad rely on infrastructure whose mathematical exposition is not yet
stable enough to import from the textbook, and their blueprint
treatment will ship once that stabilizes.

\section*{Relationship to the companion textbook}

Chapters~\ref{chap:primitives} through \ref{ch:separations} of this
blueprint are imported verbatim from the companion textbook with
\texttt{\textbackslash lean\{\}} annotations added at each theorem or
definition whose informal statement corresponds to a formal
declaration. The prose is otherwise unchanged. Figures are reused from
the textbook.

Chapter~\ref{chap:measurability} is the one chapter whose content does
not appear in the textbook, because the Borel-analytic separation and
the \texttt{WellBehavedVCMeasTarget} refinement of the fundamental
theorem were discovered during the formalization work and are outside
the textbook's scope. The prose in that chapter is therefore written
from scratch in the voice of the kernel rather than imported.

\section*{How to read this document}

A reader interested in the mathematical content can read the chapters
linearly and ignore the hyperlink annotations. A reader interested in
the formalization can click any
\texttt{\textbackslash lean\{\}}-annotated statement to open the
corresponding Lean declaration page in the doc-gen4 API reference. A
reader interested in the dependency structure of the theorems can
consult the automatically generated dependency graph at
\texttt{dep\_graph\_document.html}, which renders the
\texttt{\textbackslash uses\{\}} edges declared throughout the
document.
